\textbf{Phần 1. Trắc nghiệm lựa chọn}
\setcounterpageref{ex}{0}

\begin{ex} %[3D1N06]
Kết quả của phép tính $4 \times 6$ là
\choice
{$20$}
{\True $24$}
{$28$}
{$32$}
\loigiai{
$4 \times 6 = 24$.\\
\textbf{Đáp án của câu hỏi là B}
}
\end{ex}

\begin{ex} %[3D1N06]
Kết quả của phép tính $28 : 4$ là
\choice
{$6$}
{$8$}
{\True $7$}
{$9$}
\loigiai{
$28 : 4 = 7$.\\
\textbf{Đáp án của câu hỏi là C}
}
\end{ex}

\begin{ex} %[3D1H06]
Kết quả của phép tính $4 \times 8$ là
\choice
{$30$}
{\True $32$}
{$36$}
{$28$}
\loigiai{
$4 \times 8 = 32$.\\
\textbf{Đáp án của câu hỏi là B}
}
\end{ex}

\begin{ex} %[3D1H06]
Kết quả của phép tính $36 : 4$ là
\choice
{$8$}
{$7$}
{$6$}
{\True $9$}
\loigiai{
$36 : 4 = 9$.\\
\textbf{Đáp án của câu hỏi là D}
}
\end{ex}

\begin{ex} %[3D1H06]
Điền số thích hợp vào chỗ chấm: $4 \times \dots = 20$.
\choice
{\True $5$}
{$4$}
{$6$}
{$7$}
\loigiai{
$20 : 4 = 5$.\\
\textbf{Đáp án của câu hỏi là A}
}
\end{ex}

\begin{ex} %[3D1H06]
Điền số thích hợp vào chỗ chấm: $\dots : 4 = 4$.
\choice
{$12$}
{$20$}
{$8$}
{\True $16$}
\loigiai{
$4 \times 4 = 16$.\\
\textbf{Đáp án của câu hỏi là D}
}
\end{ex}

\begin{ex} %[3D1V06]
Tính giá trị biểu thức: $4 \times 7 + 12$.
\choice
{$38$}
{\True $40$}
{$42$}
{$44$}
\loigiai{
$4 \times 7 + 12 = 28 + 12 = 40$.\\
\textbf{Đáp án của câu hỏi là B}
}
\end{ex}

\begin{ex} %[3D1V06]
Tính giá trị biểu thức: $40 : 4 - 3$.
\choice
{$6$}
{$8$}
{\True $7$}
{$5$}
\loigiai{
$40 : 4 - 3 = 10 - 3 = 7$.\\
\textbf{Đáp án của câu hỏi là C}
}
\end{ex}

\begin{ex} %[3D1V06]
Một con bò có 4 chân. Hỏi 8 con bò có tất cả bao nhiêu cái chân?
\choice
{\True $32$ cái chân}
{$28$ cái chân}
{$36$ cái chân}
{$24$ cái chân}
\loigiai{
$4 \times 8 = 32$ cái chân.\\
\textbf{Đáp án của câu hỏi là A}
}
\end{ex}

\begin{ex} %[3D1V06]
Có 32 học sinh xếp đều thành 4 hàng. Mỗi hàng có số học sinh là
\choice
{$7$ học sinh}
{$9$ học sinh}
{$6$ học sinh}
{\True $8$ học sinh}
\loigiai{
$32 : 4 = 8$ học sinh.\\
\textbf{Đáp án của câu hỏi là D}
}
\end{ex}

\begin{ex} %[3D1V06]
Một hình vuông có cạnh dài $6\text{ cm}$. Chu vi hình vuông đó là
\choice
{$20\text{ cm}$}
{$28\text{ cm}$}
{\True $24\text{ cm}$}
{$18\text{ cm}$}
\loigiai{
$6 \times 4 = 24\text{ cm}$.\\
\textbf{Đáp án của câu hỏi là C}
}
\end{ex}

\begin{ex} %[3D1C06]
Có 24 quả cam chia đều vào các đĩa, mỗi đĩa 4 quả cam. Người ta đã biếu bà 3 đĩa cam. Hỏi còn lại bao nhiêu quả cam?
\choice
{\True $12$ quả}
{$10$ quả}
{$14$ quả}
{$8$ quả}
\loigiai{
Tổng số đĩa cam: $24 : 4 = 6$ đĩa.\\
Còn lại: $6 - 3 = 3$ đĩa $\Rightarrow 3 \times 4 = 12$ quả cam.\\
\textbf{Đáp án của câu hỏi là A}
}
\end{ex}

\textbf{Phần 2. Trắc nghiệm đúng sai}
\setcounterpageref{ex}{0}

\begin{ex} %[3D1H06]
Xét tính đúng sai của các phép tính trong bảng nhân 4, bảng chia 4:
\choiceTF
{\True $4 \times 5 = 20$.}
{\True $16 : 4 = 4$.}
{$4 \times 9 = 32$.}
{\True $4 \times 9 = 36$.}
\loigiai{
\begin{itemize}
    \item $4 \times 5 = 20$: Đúng.
    \item $16 : 4 = 4$: Đúng.
    \item $4 \times 9 = 36 \neq 32$. Ý 3 Sai, Ý 4 Đúng.
\end{itemize}
\begin{itemize}
    \item a) \textbf{Đúng}.
    \item b) \textbf{Đúng}.
    \item c) \textbf{Sai}.
    \item d) \textbf{Đúng}.
\end{itemize}
}
\end{ex}

\begin{ex} %[3D1V06]
Một sợi dây dài $36\text{ cm}$ được cắt thành 4 đoạn bằng nhau. Xét tính đúng sai:
\choiceTF
{\True Mỗi đoạn dây dài $9\text{ cm}$.}
{\True Hai đoạn dây dài $18\text{ cm}$.}
{Ba đoạn dây dài $24\text{ cm}$.}
{\True Ba đoạn dây dài $27\text{ cm}$.}
\loigiai{
\begin{itemize}
    \item Mỗi đoạn: $36 : 4 = 9\text{ cm}$. Đúng.
    \item 2 đoạn: $9 \times 2 = 18\text{ cm}$. Đúng.
    \item 3 đoạn: $9 \times 3 = 27\text{ cm} \neq 24\text{ cm}$. Ý 3 Sai, Ý 4 Đúng.
\end{itemize}
\begin{itemize}
    \item a) \textbf{Đúng}.
    \item b) \textbf{Đúng}.
    \item c) \textbf{Sai}.
    \item d) \textbf{Đúng}.
\end{itemize}
}
\end{ex}

\textbf{Phần 3. Trả lời ngắn}
\setcounterpageref{ex}{0}

\begin{ex} %[3D1H06]
Tính tích: $4 \times 4$.
\shortans{16}
\loigiai{
\textbf{Đáp số:} $16$\\
$4 \times 4 = 16$.
}
\end{ex}

\begin{ex} %[3D1H06]
Tính thương: $20 : 4$.
\shortans{5}
\loigiai{
\textbf{Đáp số:} $5$\\
$20 : 4 = 5$.
}
\end{ex}

\begin{ex} %[3D1V06]
Tính giá trị biểu thức: $4 \times 9 - 16$.
\shortans{20}
\loigiai{
\textbf{Đáp số:} $20$\\
$36 - 16 = 20$.
}
\end{ex}

\begin{ex} %[3D1V06]
Một chiếc xe ô tô có 4 bánh xe. Hỏi 7 chiếc xe ô tô như thế có bao nhiêu bánh xe?
\shortans{28}
\loigiai{
\textbf{Đáp số:} $28$\\
$4 \times 7 = 28$ bánh xe.
}
\end{ex}

\begin{ex} %[3D1V06]
Một hình vuông có chu vi $36\text{ cm}$. Tính độ dài cạnh của hình vuông đó theo xăng-ti-mét.
\shortans{9}
\loigiai{
\textbf{Đáp số:} $9$\\
$36 : 4 = 9\text{ cm}$.
}
\end{ex}

\begin{ex} %[3D1C06]
Tìm $x$ biết: $x \times 4 = 40 - 8$.
\shortans{8}
\loigiai{
\textbf{Đáp số:} $8$\\
$x \times 4 = 32 \Rightarrow x = 32 : 4 = 8$.
}
\end{ex}
